\section{Protocol Walkthrough}

This section gives a concrete end-to-end walkthrough of one Ghostext transmission. The purpose is not to provide a benchmark result, but to make the protocol mechanics and synchronization assumptions explicit in a way that can be audited and re-implemented.

\subsection{Inputs and Shared State}

Alice starts from three user-facing inputs: plaintext message, passphrase, and prompt text. Bob must share the same passphrase and prompt text. In addition, both sides must share the same protocol-sensitive runtime configuration and backend identity. In Ghostext, these parameters are not implicit; they are serialized into structured configuration and represented by a fingerprint in the packet header.

The walkthrough below assumes a fixed backend model family, fixed seed, fixed candidate policy, fixed quantization total frequency, and fixed segment token budgets.

\subsection{Step A: Plaintext to Authenticated Packet}

Alice first encodes plaintext string into UTF-8 bytes. Let this byte string be $M$.

Ghostext then derives a master key from \texttt{passphrase + salt} using scrypt and uses HKDF to split this into header and body subkeys. It encrypts body bytes with ChaCha20-Poly1305 under body key and body nonce, then builds an internal header containing protocol version, crypto identifiers, body ciphertext length, and configuration fingerprint. This internal header is itself sealed with ChaCha20-Poly1305 under header key and header nonce.

At this point, Alice obtains the opaque packet byte string:
\begin{equation}
\texttt{packet} = \texttt{salt} \parallel \texttt{header\_nonce} \parallel \texttt{sealed\_header} \parallel \texttt{body\_nonce} \parallel \texttt{body\_ct}.
\end{equation}

Ghostext then splits this packet into bootstrap bytes and body bytes. The bootstrap size is deterministic from crypto configuration.

\subsection{Step B: Bootstrap Segment Embedding}

Alice initializes an interval encoder for bootstrap bytes. At each token step $t$, Ghostext obtains backend logits for the next token conditioned on prompt and generated history. Candidate selection and quantization produce a discrete distribution with integer cumulative boundaries.

If the step is encodable, interval narrowing chooses a token whose subinterval contains the target interval point. If the step is not encodable under policy rules, Ghostext emits the top token deterministically.

This process repeats until the bootstrap interval resolves to width one. If token budget is exhausted before resolution, encoding fails.

\subsection{Step C: Body Segment Embedding}

The exact same mechanism is used for the body bytes, with a separate token budget. This separation is important because Bob can parse body length only after decoding and authenticating the bootstrap header.

In practice, this two-stage design gives cleaner failure behavior. Header-stage failures and body-stage failures are distinguishable, which later helps diagnosis and reproducibility reporting.

\subsection{Step D: Optional Natural Tail}

After both segments resolve, Ghostext may sample a short natural tail to improve output ending quality. These tokens carry no payload and are not required for decoding. This policy improves surface fluency without adding decode ambiguity, because decode consumes only the packet-resolving prefix.

\subsection{Step E: Bob's Decode Replay}

Bob tokenizes observed text under the same backend and prompt setup. He starts with bootstrap decoding. At each step, he reproduces candidate and quantization logic, then verifies that observed token is consistent with current candidate set. If inconsistent, decode fails with synchronization error.

When bootstrap bytes resolve, Bob decrypts and verifies the sealed internal header. Fingerprint mismatch immediately fails decoding. If successful, Bob learns expected body length and then decodes body bytes using the same replay logic.

Finally, Bob reconstructs full packet, verifies authenticated encryption, and recovers plaintext bytes.

\subsection{Why This Walkthrough Matters}

The key value of this walkthrough is that every protocol step corresponds to concrete implementation modules and explicit failure classes. This property is crucial for a security paper under active development: it allows collaborators to debate and improve specific transitions rather than abstract claims.

The walkthrough also clarifies where future optimization can be added safely. Distribution optimization, for example, can be inserted before quantization as long as deterministic replay, fingerprint commitments, and failure semantics remain intact.

\subsection{Compact Pseudocode View}

For clarity, the encode/decode core can be summarized as:

\begin{quote}
Encode:
packetize $\rightarrow$ split(header, body) $\rightarrow$ interval-embed(header) $\rightarrow$ interval-embed(body) $\rightarrow$ optional-tail.

Decode:
tokenize $\rightarrow$ interval-decode(header) $\rightarrow$ authenticate-header $\rightarrow$ interval-decode(body) $\rightarrow$ authenticate-body $\rightarrow$ plaintext.
\end{quote}

This compact view is intentionally simple. The implementation details that make it robust are deterministic candidate construction, integer quantization, explicit synchronization checks, and fail-closed exception handling.

\subsection{Bounded Guarantees From This Walkthrough}

Under matched configuration and passive observation assumptions, the walkthrough justifies two bounded guarantees in the current draft. First, decode correctness is deterministic when no failure condition is triggered. Second, any critical mismatch is surfaced as explicit failure rather than silent corruption.

These are strong engineering properties, but they are not complete steganographic security results. Distributional indistinguishability and detector resistance still require dedicated post-pause experiments.
